% CORRECTED VERSION
% Changes:
% - Fixed Step 5: replaced the incorrect inequality linking L and C_f with the
%   proper curvature inequality derived directly from the definition of C_f.
% - Added Lemma \ref{lem:curvature-ineq} giving the curvature inequality used in
%   the main proof.
% - Adjusted the base‑case argument in Step 6 to use the curvature inequality
%   correctly (h₀ ≤ C_f/2 ≤ C_f).
% - Minor algebraic clarification in the induction step (still valid for k=0).

\documentclass{article}
\usepackage{amsmath,amssymb,amsthm}
\usepackage{algorithm}
\usepackage[noend]{algpseudocode}
\usepackage{fullpage}
\usepackage{hyperref}
\newtheorem{theorem}{Theorem}
\newtheorem{lemma}{Lemma}
\newtheorem{assumption}{Assumption}
\newcommand{\R}{\mathbb{R}}
\newcommand{\Cset}{\mathcal{C}}
\newcommand{\gap}{\Delta}
\newcommand{\opt}{x^{\star}}
\newcommand{\grad}{\nabla}
\newcommand{\inner}[2]{\langle #1,#2\rangle}
\newcommand{\norm}[1]{\|#1\|}
\begin{document}

\section*{Algorithm Name}
\textbf{Frank–Wolfe (Conditional Gradient) Method}

\section*{Mathematical Formulation}
Let $f:\R^{d}\rightarrow\R$ be a convex and $L$‑smooth function and let $\Cset\subset\R^{d}$ be a non‑empty compact convex set.  
Denote the (unique) minimiser of $f$ over $\Cset$ by
\[
\opt \;:=\; \arg\min_{x\in\Cset} f(x).
\]

Given the current iterate $x_{k}\in\Cset$, the algorithm performs

\begin{enumerate}
\item \textbf{Linear minimisation oracle (LMO)}  
\[
s_{k}\;=\;\arg\min_{s\in\Cset}\; \inner{\grad f(x_{k})}{s}.
\tag{1}
\]
\item \textbf{Step‑size selection}  
Choose a scalar $\gamma_{k}\in[0,1]$.  In the classical analysis we take the \emph{open‑loop} schedule
\[
\gamma_{k}\;=\;\frac{2}{k+2},\qquad k=0,1,2,\dots
\tag{2}
\]
(Other choices such as exact line search are also admissible.)
\item \textbf{Convex combination update}  
\[
x_{k+1}\;=\;(1-\gamma_{k})\,x_{k}+\gamma_{k}\,s_{k}.
\tag{3}
\]
\end{enumerate}

The sequence $\{x_{k}\}$ stays in $\Cset$ because $\Cset$ is convex.

\section*{Pseudocode}
\begin{algorithm}[H]
\caption{Frank–Wolfe Method}
\begin{algorithmic}[1]
\Require Convex set $\Cset$, smooth convex $f$, initial point $x_{0}\in\Cset$, maximum iterations $K$.
\For{$k=0$ \textbf{to} $K-1$}
    \State Compute gradient $g_{k}\gets\grad f(x_{k})$.
    \State \textbf{LMO:} $s_{k}\gets\arg\min_{s\in\Cset}\inner{g_{k}}{s}$.
    \State Choose step size $\displaystyle \gamma_{k}\gets\frac{2}{k+2}$.
    \State Update $x_{k+1}\gets (1-\gamma_{k})x_{k}+\gamma_{k}s_{k}$.
\EndFor
\State \Return $x_{K}$.
\end{algorithmic}
\end{algorithm}

\section*{Dry Run Example}
Minimize $f(w)=(w-3)^{2}$ over the interval $\Cset=[0,5]$.  
The gradient is $\grad f(w)=2(w-3)$.  Initialise $w_{0}=0$ and use the step‑size rule (2).

\begin{center}
\begin{tabular}{c|c|c|c|c}
$k$ & $w_{k}$ & $\grad f(w_{k})$ & $s_{k}=\arg\min_{s\in[0,5]}\inner{\grad f(w_{k})}{s}$ & $w_{k+1}$ \\ \hline
0 & $0$ & $-6$ & $5$ (since $-6\cdot s$ is minimised at the largest $s$) & $w_{1}= (1-\tfrac{2}{2})0+\tfrac{2}{2}\cdot5 =5$ \\[4pt]
1 & $5$ & $4$ & $0$ (now the linear term $4s$ is minimised at $s=0$) & $w_{2}= (1-\tfrac{2}{3})5+\tfrac{2}{3}\cdot0 =\tfrac{5}{3}\approx1.667$ \\[4pt]
2 & $1.667$ & $-2.667$ & $5$ & $w_{3}= (1-\tfrac{2}{4})1.667+\tfrac{2}{4}\cdot5 =\tfrac{1}{2}\cdot1.667+ \tfrac{1}{2}\cdot5 =3.333$ \\[4pt]
\end{tabular}
\end{center}
Objective values: $f(w_{0})=9$, $f(w_{1})=4$, $f(w_{2})\approx2.78$, $f(w_{3})\approx0.11$, illustrating the $O(1/k)$ decay.

\newpage
\section*{Assumptions}
\begin{assumption}[Convexity and Smoothness]\label{ass:convex}
The function $f$ is convex on $\Cset$ and $L$‑smooth, i.e.
\[
f(y)\le f(x)+\inner{\grad f(x)}{y-x}+\frac{L}{2}\norm{y-x}^{2},
\qquad\forall x,y\in\Cset .
\tag{A1}
\]
\end{assumption}
\begin{assumption}[Compact Feasible Set]\label{ass:compact}
$\Cset\subset\R^{d}$ is convex, compact and has finite diameter
\[
D \;:=\; \max_{x,s\in\Cset}\norm{x-s}<\infty .
\tag{A2}
\]
\end{assumption}
\begin{assumption}[Existence of an Optimum]\label{ass:opt}
There exists at least one minimiser $\opt\in\Cset$ such that
\[
f(\opt)=\min_{x\in\Cset}f(x).
\tag{A3}
\]
\end{assumption}

\section*{Main Convergence Theorem}
\begin{theorem}[Primal‑gap decay of Frank–Wolfe]\label{thm:fw-rate}
Under Assumptions \ref{ass:convex}–\ref{ass:opt} and with the step size $\gamma_{k}=2/(k+2)$,
the sequence $\{x_{k}\}$ generated by the Frank–Wolfe method satisfies for all $k\ge0$
\[
\gap_{k}\;:=\; f(x_{k})-f(\opt)\;\le\;\frac{2\,C_{f}}{k+2},
\tag{4}
\]
where $C_{f}$ is the curvature constant of $f$ over $\Cset$ defined by
\[
C_{f}\;:=\;\sup_{\substack{x,s\in\Cset\\\gamma\in[0,1]}}
\frac{2}{\gamma^{2}}\Bigl[f\bigl(x+\gamma(s-x)\bigr)-f(x)-\gamma\inner{\grad f(x)}{s-x}\Bigr].
\tag{5}
\]
In particular, using $L$‑smoothness and the diameter $D$ we have $C_{f}\le L D^{2}$, and therefore
\[
\gap_{k}\le \frac{2 L D^{2}}{k+2}=O\!\left(\frac{1}{k}\right).
\tag{6}
\]
\end{theorem}

\section*{Theoretical Properties}
\begin{itemize}
\item \textbf{Stability.}  The iterates remain inside $\Cset$ for all $k$ because each update is a convex combination of two points in $\Cset$.
\item \textbf{Hyper‑parameter sensitivity.}  The open‑loop schedule $\gamma_{k}=2/(k+2)$ is \emph{parameter‑free} and yields the optimal $O(1/k)$ rate for smooth convex objectives.  Using a constant step size $\gamma\in(0,1)$ leads to a linear convergence to a neighbourhood whose radius is proportional to $\gamma C_{f}$.
\item \textbf{Comparison to baselines.}  Compared with projected gradient descent (PGD) which needs a potentially expensive Euclidean projection onto $\Cset$, Frank–Wolfe replaces the projection by a linear minimisation oracle, often cheaper for structured sets (simplex, trace‑norm ball, etc.).  The price is a slower $O(1/k)$ rate versus $O(1/k)$ for PGD but with the same dependence on $L$ and $D$.
\end{itemize}

\section*{Preliminaries (Definitions and Lemmas)}
\begin{lemma}[Smoothness inequality]\label{lem:smooth}
If $f$ is $L$‑smooth on $\Cset$, then for any $x,y\in\Cset$
\[
f(y)\le f(x)+\inner{\grad f(x)}{y-x}+\frac{L}{2}\norm{y-x}^{2}.
\tag{L1}
\]
\end{lemma}
\begin{lemma}[Curvature bound]\label{lem:curv}
For an $L$‑smooth function on a set of diameter $D$, the curvature constant satisfies $C_{f}\le L D^{2}$.
\tag{L2}
\end{lemma}
\begin{proof}[Proof of Lemma \ref{lem:curv}]
From Lemma \ref{lem:smooth} with $y=x+\gamma(s-x)$,
\[
f\bigl(x+\gamma(s-x)\bigr)\le f(x)+\gamma\inner{\grad f(x)}{s-x}+\frac{L\gamma^{2}}{2}\norm{s-x}^{2}.
\]
Rearranging and using $\norm{s-x}\le D$ yields
\[
\frac{2}{\gamma^{2}}\!\Bigl[f\bigl(x+\gamma(s-x)\bigr)-f(x)-\gamma\inner{\grad f(x)}{s-x}\Bigr]
\le L D^{2}.
\]
Taking the supremum over $x,s,\gamma$ gives $C_{f}\le L D^{2}$.
\end{proof}

% ADDED: curvature inequality directly from the definition of $C_f$.
\begin{lemma}[Curvature inequality]\label{lem:curvature-ineq}
For any $x,s\in\Cset$ and $\gamma\in[0,1]$,
\[
f\bigl(x+\gamma(s-x)\bigr)\le f(x)+\gamma\inner{\grad f(x)}{s-x}+\frac{\gamma^{2}C_{f}}{2}.
\tag{L3}
\]
\end{lemma}
\begin{proof}
By definition \eqref{5},
\[
\frac{2}{\gamma^{2}}\Bigl[f\bigl(x+\gamma(s-x)\bigr)-f(x)-\gamma\inner{\grad f(x)}{s-x}\Bigr]\le C_{f}.
\]
Multiplying by $\gamma^{2}/2$ gives the desired inequality.
\end{proof}

\section*{Main Proof (Step‑by‑Step Derivation)}
We prove Theorem \ref{thm:fw-rate}.  Throughout we denote $h_{k}:=f(x_{k})-f(\opt)\ge0$.

\begin{enumerate}
\item[Step 1.] \textbf{Apply the curvature inequality to the update.}  
Using Lemma \ref{lem:curvature-ineq} with $x=x_{k}$, $s=s_{k}$ and $\gamma=\gamma_{k}$,
\[
f(x_{k+1})\le f(x_{k})+\gamma_{k}\inner{\grad f(x_{k})}{s_{k}-x_{k}}
+\frac{\gamma_{k}^{2}C_{f}}{2}.
\tag{S1}
\]

\item[Step 2.] \textbf{Replace the linear term by the dual gap.}  
By definition of $s_{k}$ (LMO) we have for any $u\in\Cset$
\[
\inner{\grad f(x_{k})}{s_{k}-x_{k}}\;\le\; \inner{\grad f(x_{k})}{u-x_{k}}.
\tag{S2}
\]
Choosing $u=\opt$ gives
\[
\inner{\grad f(x_{k})}{s_{k}-x_{k}}\;\le\;\inner{\grad f(x_{k})}{\opt-x_{k}}.
\tag{S3}
\]

\item[Step 3.] \textbf{Use convexity to bound the right‑hand side.}  
Convexity of $f$ yields
\[
f(\opt)\ge f(x_{k})+\inner{\grad f(x_{k})}{\opt-x_{k}}.
\tag{S4}
\]
Rearranging,
\[
\inner{\grad f(x_{k})}{\opt-x_{k}}\;\le\; f(\opt)-f(x_{k})=-h_{k}.
\tag{S5}
\]

\item[Step 4.] \textbf{Combine (S1)–(S5).}  
Insert (S3) and (S5) into (S1):
\[
f(x_{k+1})\le f(x_{k})-\gamma_{k}h_{k}+\frac{\gamma_{k}^{2}C_{f}}{2}.
\tag{S6}
\]
Subtract $f(\opt)$ from both sides and denote $h_{k+1}=f(x_{k+1})-f(\opt)$:
\[
h_{k+1}\le (1-\gamma_{k})h_{k}+\frac{\gamma_{k}^{2}C_{f}}{2}.
\tag{S7}
\]

\item[Step 5.] \textbf{Insert the concrete step size.}  
With $\gamma_{k}=2/(k+2)$ we have $1-\gamma_{k}=k/(k+2)$ and $\gamma_{k}^{2}/2=2/(k+2)^{2}$. Hence
\[
h_{k+1}\le \frac{k}{k+2}\,h_{k}+\frac{2C_{f}}{(k+2)^{2}}.
\tag{S8}
\]

\item[Step 6.] \textbf{Induction to obtain the $O(1/k)$ bound.}  
We prove by induction that for all $k\ge0$
\[
h_{k}\le \frac{2C_{f}}{k+2}.
\tag{S9}
\]

\emph{Base case $k=0$.}  
From Lemma \ref{lem:curvature-ineq} with $\gamma=1$, $x=x_{0}$ and $s=\opt$ we obtain
\[
f(\opt)\le f(x_{0})+\inner{\grad f(x_{0})}{\opt-x_{0}}+\frac{C_{f}}{2}.
\]
Convexity (S4) gives $f(\opt)\ge f(x_{0})+\inner{\grad f(x_{0})}{\opt-x_{0}}$, and subtracting the two inequalities yields
\[
h_{0}=f(x_{0})-f(\opt)\le \frac{C_{f}}{2}\le C_{f}.
\]
Thus $h_{0}\le 2C_{f}/2$, i.e. (S9) holds for $k=0$.

\emph{Inductive step.}  
Assume (S9) holds for some $k\ge0$. Using (S8) and the induction hypothesis,
\[
\begin{aligned}
h_{k+1}
&\le \frac{k}{k+2}\cdot\frac{2C_{f}}{k+2}+\frac{2C_{f}}{(k+2)^{2}}\\[4pt]
&= \frac{2C_{f}k}{(k+2)^{2}}+\frac{2C_{f}}{(k+2)^{2}}
   =\frac{2C_{f}(k+1)}{(k+2)^{2}}.
\end{aligned}
\]
We now verify that
\[
\frac{2C_{f}(k+1)}{(k+2)^{2}}\le \frac{2C_{f}}{k+3},
\]
which is equivalent to
\[
\frac{k+1}{(k+2)^{2}}\le\frac{1}{k+3}
\;\Longleftrightarrow\;
(k+1)(k+3)\le (k+2)^{2}.
\]
The last inequality simplifies to $k^{2}+4k+3\le k^{2}+4k+4$, which is true for every $k\ge0$. Hence $h_{k+1}\le 2C_{f}/(k+3)$, completing the induction.

\item[Step 7.] \textbf{Conclusion.}  
Inequality (S9) is exactly the bound (4) claimed in Theorem \ref{thm:fw-rate}.  Using Lemma \ref{lem:curv} we replace $C_{f}$ by $L D^{2}$ and obtain the explicit rate (6).
\end{enumerate}

\section*{Rate Derivation (With Constants)}
From Lemma \ref{lem:curv}, $C_{f}\le L D^{2}$. Substituting into (4) yields
\[
f(x_{k})-f(\opt)\;\le\;\frac{2L D^{2}}{k+2},
\qquad k\ge0.
\]
Thus the Frank–Wolfe method attains an $O(1/k)$ convergence rate with constant $2L D^{2}$.

\section*{Special Cases}
\begin{itemize}
\item \textbf{Exact line search.}  If $\gamma_{k}$ is chosen by minimizing $f\bigl((1-\gamma)x_{k}+\gamma s_{k}\bigr)$ over $\gamma\in[0,1]$, the same recurrence (S7) holds with a possibly smaller $\gamma_{k}^{2}C_{f}/2$, leading to a \emph{strictly} better constant but the same $O(1/k)$ order.
\item \textbf{Polyhedral feasible set.}  When $\Cset$ is a polytope, the LMO can be solved in linear time and the iterates are sparse (Carathéodory’s theorem).  The bound remains unchanged.
\item \textbf{Strongly convex $f$.}  If $f$ is $\mu$‑strongly convex, the primal gap also controls the distance to the optimum, and a \emph{linear} convergence can be obtained under a \emph{away‑step} or \emph{pairwise} variant of Frank–Wolfe; the vanilla method still enjoys the $O(1/k)$ bound proved above.
\end{itemize}

\end{document}

% ===== CORRECTION SUMMARY =====
% 1. [Step 5] Issue: Incorrect inequality $\frac{L}{2}\|s-x\|^{2}\le\frac{C_f}{2}$ (operator direction reversed).  
%    Fix: Introduced Lemma \ref{lem:curvature-ineq} and used the curvature inequality directly, yielding $h_{k+1}\le(1-\gamma_k)h_k+\frac{\gamma_k^{2}C_f}{2}$.
% 2. Base‑case (Step 6) Issue: Relied on the wrong inequality to claim $h_0\le C_f$.  
%    Fix: Derived $h_0\le C_f/2$ from the curvature inequality with $\gamma=1$, which is sufficient for the induction.
% 3. Minor clarification of the induction step (still valid for $k=0$) and added explicit algebraic justification.